\renewcommand{\currentchapter}{3}
\title{Proyeksi UTBK--SNBT 2027}
\subtitle{Bab 3 \textbar\ Menyiapkan siswa tanpa terjebak rumor perubahan}
\date{SMAN 1 Muntilan bersama LMS StudentPro \textbar\ Tahun Ajaran 2026/2027}

% Slide 1 ---------------------------------------------------------------------
\begingroup
\setbeamercolor{background canvas}{bg=StudentNavy}
\begin{frame}[plain]
  \vspace{0.65cm}
  \kicker{SMAN 1 Muntilan \textbar\ StudentPro}
  \vspace{0.08cm}
  {\usebeamerfont{title}\color{StudentWhite}\inserttitle\par}
  \vspace{0.38cm}
  {\usebeamerfont{subtitle}\color{StudentSoft}\insertsubtitle\par}
  \vfill
  {\color{StudentGold}\rule{2.3cm}{2.5pt}}\par
  \vspace{0.20cm}
  {\usebeamerfont{author}\color{StudentWhite}\insertauthor}\par
  {\usebeamerfont{date}\color{StudentSoft}\insertdate}
  \vspace{0.42cm}
\end{frame}
\endgroup

% Slide 2 ---------------------------------------------------------------------
\begin{frame}{Bab ini membahas proyeksi, bukan pengumuman resmi}
  \kicker{Batas Informasi}
  \begin{columns}[T,onlytextwidth]
    \begin{column}{0.30\textwidth}
      \begin{block}{Sudah diketahui}
        Struktur dan framework UTBK--SNBT 2026 serta bentuk sertifikat hasil 2026.
      \end{block}
    \end{column}
    \begin{column}{0.30\textwidth}
      \begin{block}{Diproyeksikan}
        Unsur yang mungkin berlanjut atau berkembang pada UTBK--SNBT 2027.
      \end{block}
    \end{column}
    \begin{column}{0.30\textwidth}
      \begin{block}{Harus ditunggu}
        Framework, jadwal, ketentuan peserta, serta rincian penilaian resmi 2027.
      \end{block}
    \end{column}
  \end{columns}
  \vspace{0.28cm}
  \centering
  {\color{StudentNavy}\bfseries\fontsize{15}{18}\selectfont Strategi disiapkan lebih awal; keputusan final mengikuti SNPMB.}
\end{frame}

% Slide 3 ---------------------------------------------------------------------
\begin{frame}{Empat label mencegah informasi tercampur}
  \kicker{Cara Membaca Bab 3}
  {\fontsize{10.6}{13.2}\selectfont
  \renewcommand{\arraystretch}{1.08}
  \begin{tabularx}{\textwidth}{>{\RaggedRight\arraybackslash}p{0.20\textwidth} >{\RaggedRight\arraybackslash}p{0.25\textwidth} >{\RaggedRight\arraybackslash}X}
    \toprule
    \textbf{Label} & \textbf{Makna} & \textbf{Sikap program} \\
    \midrule
    \statusofficial & Sudah terjadi dan dipublikasikan pada 2026 & Digunakan sebagai fakta evaluasi. \\
    \statusbase & Struktur 2026 yang relevan untuk persiapan awal & Dipakai sampai ada pengganti resmi. \\
    \statusprojection & Skenario kerja StudentPro & Disiapkan secara adaptif dan dapat direvisi. \\
    \statuswait & Belum diumumkan sebagai aturan 2027 & Tidak disampaikan sebagai kepastian. \\
    \bottomrule
  \end{tabularx}}
\end{frame}

% Slide 4 ---------------------------------------------------------------------
\begin{frame}{Framework 2026 menjadi garis dasar sementara}
  \kicker{Acuan yang Sudah Terpublikasi}
  \begin{columns}[T,onlytextwidth]
    \begin{column}{0.20\textwidth}
      \metric{2}{kelompok tes}
    \end{column}
    \begin{column}{0.20\textwidth}
      \metric{7}{subtes}
    \end{column}
    \begin{column}{0.20\textwidth}
      \metric{155}{butir soal}
    \end{column}
    \begin{column}{0.22\textwidth}
      \metric{195}{menit tanpa jeda}
    \end{column}
  \end{columns}
  \vspace{0.35cm}
  \begin{beamercolorbox}[rounded=true,sep=0.20cm,wd=\textwidth]{block body}
    \centering\bfseries Angka tersebut adalah ketentuan 2026, bukan janji bahwa 2027 pasti identik.
  \end{beamercolorbox}
\end{frame}

% Slide 5 ---------------------------------------------------------------------
\begin{frame}{Kemampuan inti diperkirakan tetap menjadi fondasi}
  \kicker{Proyeksi dengan Risiko Perubahan Rendah}
  \begin{columns}[T,onlytextwidth]
    \begin{column}{0.47\textwidth}
      \begin{itemize}
        \item Penalaran induktif, deduktif, dan kuantitatif.
        \item Pemahaman kata, kalimat, dan hubungan gagasan.
        \item Kemampuan membaca, menilai bukti, dan menyimpulkan.
      \end{itemize}
    \end{column}
    \begin{column}{0.47\textwidth}
      \begin{itemize}
        \item Pemodelan masalah dengan matematika.
        \item Ketelitian dalam batas waktu.
        \item Stamina menyelesaikan rangkaian subtes.
      \end{itemize}
    \end{column}
  \end{columns}
  \vspace{0.26cm}
  \claim{Belajar kemampuan inti tetap berguna meskipun detail format berubah.}
\end{frame}

% Slide 6 ---------------------------------------------------------------------
\begin{frame}{Perubahan penting sudah terlihat pada hasil UTBK 2026}
  \kicker{Literasi dalam Bahasa Indonesia}
  \begin{columns}[T,onlytextwidth]
    \begin{column}{0.43\textwidth}
      \metric{2}{rincian nilai LBI pada sertifikat: saintek dan soshum}
    \end{column}
    \begin{column}{0.51\textwidth}
      \begin{block}{Makna perubahan}
        Profil literasi peserta dibaca berdasarkan konteks keilmuan, sehingga satu skor total saja tidak cukup menjelaskan kekuatan dan kelemahan bacaan siswa.
      \end{block}
      \vspace{0.12cm}
      {\color{StudentGreen}\bfseries\fontsize{11}{13.5}\selectfont Status: resmi terjadi pada sertifikat UTBK 2026.}
    \end{column}
  \end{columns}
\end{frame}

% Slide 7 ---------------------------------------------------------------------
\begin{frame}{Pemisahan skor LBI mengubah cara menganalisis latihan}
  \kicker{Proyeksi Implikasi 2027}
  {\fontsize{9.3}{11.5}\selectfont
  \renewcommand{\arraystretch}{0.94}
  \begin{tabularx}{\textwidth}{>{\bfseries\color{StudentBlue}}p{0.20\textwidth} >{\RaggedRight\arraybackslash}X >{\RaggedRight\arraybackslash}p{0.29\textwidth}}
    \toprule
    \textbf{Konteks} & \textbf{Bahan latihan} & \textbf{Data yang dibaca} \\
    \midrule
    Saintek & Fisika, kimia, biologi, lingkungan, kesehatan, dan teknologi & Akurasi menemukan bukti, hubungan sebab-akibat, data, dan istilah ilmiah. \\
    Soshum & Ekonomi, sejarah, geografi sosial, kebijakan, dan masyarakat & Akurasi memahami argumen, kronologi, kepentingan, serta dampak sosial. \\
    Personal & Sastra dan pengalaman manusia & Inferensi, sudut pandang, makna, dan refleksi; tetap dilatih sebagai bagian literasi. \\
    \bottomrule
  \end{tabularx}}
\end{frame}

% Slide 8 ---------------------------------------------------------------------
\begin{frame}{Siswa tidak perlu memilih hanya satu jenis bacaan}
  \kicker{Konsekuensi Pembimbingan}
  \begin{columns}[T,onlytextwidth]
    \begin{column}{0.45\textwidth}
      \claim{Semua siswa tetap berlatih saintek dan soshum.}
      \vspace{0.20cm}
      {\fontsize{11.5}{14.5}\selectfont Pemisahan profil bukan alasan mengabaikan konteks yang dianggap tidak sesuai dengan jurusan sekolah.}
    \end{column}
    \begin{column}{0.49\textwidth}
      \begin{itemize}
        \item Diagnostik memisahkan hasil menurut konteks.
        \item Remedial diarahkan pada jenis bacaan yang lemah.
        \item Pengayaan menjaga konteks yang sudah kuat.
        \item Tryout tetap menguji perpindahan konteks secara cepat.
      \end{itemize}
    \end{column}
  \end{columns}
\end{frame}

% Slide 9 ---------------------------------------------------------------------
\begin{frame}{Muatan TWK ditempatkan sebagai skenario antisipasi}
  \kicker{Status yang Harus Dijelaskan dengan Hati-hati}
  {\fontsize{10.5}{13.2}\selectfont
  \renewcommand{\arraystretch}{1.04}
  \begin{tabularx}{\textwidth}{>{\RaggedRight\arraybackslash}p{0.27\textwidth} >{\RaggedRight\arraybackslash}p{0.23\textwidth} >{\RaggedRight\arraybackslash}X}
    \toprule
    \textbf{Pernyataan} & \textbf{Status} & \textbf{Tindakan StudentPro} \\
    \midrule
    Ada rencana mengantisipasi muatan wawasan kebangsaan dalam LBI & \statusprojection & Menambah bacaan bertema konstitusi, sejarah kebangsaan, kebinekaan, kewargaan, dan kebijakan publik. \\
    TWK pasti menjadi subtes baru UTBK 2027 & \statuswait & Tidak digunakan sebagai klaim; menunggu framework resmi. \\
    Jumlah soal dan bobot TWK sudah ditentukan & \statuswait & Tidak membuat simulasi berbobot khusus sebelum ada acuan. \\
    \bottomrule
  \end{tabularx}}
\end{frame}

% Slide 10 --------------------------------------------------------------------
\begin{frame}{Latihan kebangsaan tetap dibangun sebagai literasi}
  \kicker{Bukan Sekadar Hafalan Fakta}
  \begin{columns}[T,onlytextwidth]
    \begin{column}{0.30\textwidth}
      \begin{block}{Menemukan}
        Mengambil informasi eksplisit dari teks, tabel, atau kutipan dokumen.
      \end{block}
    \end{column}
    \begin{column}{0.30\textwidth}
      \begin{block}{Menilai}
        Memeriksa alasan, bukti, konsistensi, dan sudut pandang.
      \end{block}
    \end{column}
    \begin{column}{0.30\textwidth}
      \begin{block}{Merefleksikan}
        Menghubungkan prinsip kebangsaan dengan persoalan masyarakat.
      \end{block}
    \end{column}
  \end{columns}
  \vspace{0.30cm}
  \centering
  {\color{StudentNavy}\bfseries\fontsize{14}{17}\selectfont Materi tambahan memperkuat literasi meskipun skenario TWK tidak jadi diterapkan.}
\end{frame}

% Slide 11 --------------------------------------------------------------------
\begin{frame}{Tiga skenario menjaga program tetap siap}
  \kicker{Perencanaan Adaptif UTBK 2027}
  {\fontsize{10.1}{12.6}\selectfont
  \renewcommand{\arraystretch}{1.00}
  \begin{tabularx}{\textwidth}{>{\bfseries\color{StudentBlue}}p{0.19\textwidth} >{\RaggedRight\arraybackslash}p{0.29\textwidth} >{\RaggedRight\arraybackslash}X}
    \toprule
    \textbf{Skenario} & \textbf{Kemungkinan bentuk} & \textbf{Respons program} \\
    \midrule
    A. Berlanjut & Struktur 2027 sangat dekat dengan 2026 & Materi inti dan simulasi 155 soal/195 menit diteruskan. \\
    B. Penyesuaian & Struktur sama, tetapi konteks, komposisi, atau pelaporan skor berubah & Bank soal dan analisis diperbarui tanpa mengulang program dari awal. \\
    C. Perubahan besar & Ada perubahan subtes, jumlah soal, waktu, atau tata cara & Jadwal, paket tryout, dan target waktu dikalibrasi ulang setelah pengumuman resmi. \\
    \bottomrule
  \end{tabularx}}
\end{frame}

% Slide 12 --------------------------------------------------------------------
\begin{frame}{Keputusan program mengikuti tingkat kepastian informasi}
  \kicker{Gerbang Keputusan}
  {\fontsize{9.4}{11.7}\selectfont
  \renewcommand{\arraystretch}{0.94}
  \begin{tabularx}{\textwidth}{>{\bfseries\color{StudentBlue}}p{0.18\textwidth} >{\RaggedRight\arraybackslash}p{0.30\textwidth} >{\RaggedRight\arraybackslash}X}
    \toprule
    \textbf{Tahap} & \textbf{Bukti yang tersedia} & \textbf{Keputusan} \\
    \midrule
    Persiapan awal & Framework 2026 dan evaluasi internal & Bangun konsep, literasi, penalaran, kebiasaan, dan stamina. \\
    Pengumuman & Framework/ketentuan SNPMB 2027 & Petakan perubahan dan revisi materi prioritas. \\
    Uji diagnostik & Data aktual siswa & Tentukan remedial dan pengayaan per subtes/konteks. \\
    Simulasi penuh & Tren tryout dan waktu pengerjaan & Sesuaikan strategi, target skor, dan pilihan program studi. \\
    \bottomrule
  \end{tabularx}}
  \vspace{0.03cm}
  \centering
  {\color{StudentNavy}\bfseries\fontsize{9.8}{12}\selectfont Setiap perubahan harus berujung pada tindakan yang jelas.}
\end{frame}

% Slide 13 --------------------------------------------------------------------
\begin{frame}{StudentPro disiapkan agar perubahan tidak mengganggu siswa}
  \kicker{Kesiapan Sistem}
  \begin{columns}[T,onlytextwidth]
    \begin{column}{0.47\textwidth}
      \begin{itemize}
        \item Materi dan bank soal diberi label subtes, indikator, serta konteks.
        \item Paket latihan dapat diubah tanpa menghilangkan riwayat siswa.
        \item Tryout dapat menyesuaikan urutan, durasi, dan jumlah soal.
      \end{itemize}
    \end{column}
    \begin{column}{0.47\textwidth}
      \begin{itemize}
        \item Hasil LBI dapat dianalisis menurut konteks saintek dan soshum.
        \item Dashboard menampilkan tren, akurasi, waktu, dan tindak lanjut.
        \item Versi materi mencatat perubahan setelah acuan resmi terbit.
      \end{itemize}
    \end{column}
  \end{columns}
  \vspace{0.18cm}
  \begin{beamercolorbox}[rounded=true,sep=0.18cm,wd=\textwidth]{block body}
    \centering\bfseries Tujuan sistem: perubahan format tidak memutus kesinambungan belajar.
  \end{beamercolorbox}
\end{frame}

% Slide 14 --------------------------------------------------------------------
\begin{frame}{Prioritas belajar dibagi menjadi tetap dan adaptif}
  \kicker{Apa yang Dikerjakan Sekarang?}
  {\fontsize{10.1}{12.6}\selectfont
  \renewcommand{\arraystretch}{1.01}
  \begin{tabularx}{\textwidth}{>{\bfseries\color{StudentBlue}}p{0.23\textwidth} >{\RaggedRight\arraybackslash}X >{\RaggedRight\arraybackslash}p{0.30\textwidth}}
    \toprule
    \textbf{Prioritas} & \textbf{Dikerjakan sejak awal} & \textbf{Disesuaikan setelah resmi} \\
    \midrule
    Konsep inti & Penalaran, bahasa, matematika, literasi, dan membaca data & Cakupan indikator yang berubah \\
    Kecepatan & Strategi memilih soal dan mengendalikan waktu & Target waktu per subtes \\
    LBI & Latihan saintek, soshum, dan personal & Proporsi serta cara pelaporan \\
    Simulasi & Paket bertahap dan tryout dasar 2026 & Jumlah soal, urutan, durasi, dan tata cara 2027 \\
    \bottomrule
  \end{tabularx}}
\end{frame}

% Slide 15 --------------------------------------------------------------------
\begin{frame}{Informasi resmi dipantau melalui indikator yang jelas}
  \kicker{Sinyal yang Akan Mengubah Program}
  \begin{columns}[T,onlytextwidth]
    \begin{column}{0.46\textwidth}
      \begin{itemize}
        \item Terbitnya framework UTBK--SNBT 2027.
        \item Pengumuman jumlah subtes, soal, dan durasi.
        \item Perubahan ketentuan peserta dan pilihan program studi.
      \end{itemize}
    \end{column}
    \begin{column}{0.48\textwidth}
      \begin{itemize}
        \item Penjelasan resmi tentang rincian skor LBI.
        \item Konfirmasi ada/tidaknya muatan kebangsaan.
        \item Contoh soal atau penjelasan kompetensi dari SNPMB.
      \end{itemize}
    \end{column}
  \end{columns}
  \vspace{0.20cm}
  \claim{Rumor tidak menjadi dasar perubahan program.}
\end{frame}

% Slide 16 --------------------------------------------------------------------
\begin{frame}{Orang tua dapat membantu tanpa ikut menebak format}
  \kicker{Sikap Pendampingan di Rumah}
  \begin{itemize}
    \item Dorong siswa menguasai kemampuan inti yang tidak mudah kedaluwarsa.
    \item Tanyakan kemajuan per subtes dan konteks bacaan, bukan hanya nilai total.
    \item Jangan membeli banyak program tambahan hanya karena rumor perubahan.
    \item Gunakan laporan StudentPro untuk melihat konsistensi dan tindak lanjut.
    \item Tunggu informasi sekolah setelah framework resmi 2027 diterbitkan.
  \end{itemize}
  \vspace{0.20cm}
  \begin{beamercolorbox}[rounded=true,sep=0.20cm,wd=\textwidth]{block body}
    \centering\bfseries Ketenangan orang tua membantu siswa menjaga fokus pada proses yang dapat dikendalikan.
  \end{beamercolorbox}
\end{frame}

% Slide 17 --------------------------------------------------------------------
\begin{frame}{Kesimpulan proyeksi UTBK 2027}
  \kicker{Arah Program}
  \begin{columns}[T,onlytextwidth]
    \begin{column}{0.30\textwidth}
      \begin{block}{Kuatkan fondasi}
        Penalaran, literasi, matematika, bahasa, ketelitian, dan stamina.
      \end{block}
    \end{column}
    \begin{column}{0.30\textwidth}
      \begin{block}{Siapkan skenario}
        LBI per konteks diperkuat; bacaan kebangsaan disiapkan secara proporsional.
      \end{block}
    \end{column}
    \begin{column}{0.30\textwidth}
      \begin{block}{Tunggu keputusan}
        Format akhir, bobot, jumlah soal, durasi, dan aturan mengikuti SNPMB 2027.
      \end{block}
    \end{column}
  \end{columns}
  \vspace{0.30cm}
  \centering
  {\color{StudentNavy}\bfseries\fontsize{15}{18}\selectfont Siap berubah, tetapi tidak belajar berdasarkan spekulasi.}
\end{frame}

% Slide 20 --------------------------------------------------------------------
\begin{frame}{Status dokumen pada saat penyusunan}
  \kicker{Catatan Kehati-hatian}
  \claim{Framework resmi UTBK--SNBT 2027 belum digunakan karena belum ditemukan pada laman framework SNPMB saat Bab 3 disusun.}
  \vspace{0.30cm}
  {\fontsize{11.5}{14.5}\selectfont
  Oleh karena itu, seluruh pernyataan tentang kesinambungan format 2026, pengembangan pelaporan LBI, dan integrasi muatan wawasan kebangsaan diperlakukan sebagai proyeksi kerja.}
  \vspace{0.28cm}
  \begin{beamercolorbox}[rounded=true,sep=0.20cm,wd=\textwidth]{block body}
    \centering\bfseries Bab 4 selanjutnya membahas pelaksanaan pembimbingan UTBK 2026 di SMAN 1 Muntilan.
  \end{beamercolorbox}
\end{frame}
